\begin{ouverture}[probabilites]{21}{Les probabilités : a chances sur b}{Partie 5 · Modéliser et organiser}
{Issues · Équiprobabilité · Arbres et tableaux}
\objectif{Exprimer une probabilité sous la forme « a chances sur b ».}
\objectif{Comparer des probabilités en dénombrant les issues.}
\objectif{Recenser les issues d’une expérience à deux étapes.}
\prerequis{Vocabulaire du hasard et fractions.}
\end{ouverture}
\pagestyle{otMain}\essentiel\fontsize{11}{15}\selectfont
\phantomsection\label{cours-21-1}
\essSection{1}{Quantifier les chances}
\begin{definition}Quand toutes les issues d’une expérience ont les mêmes chances - dé équilibré, boules indiscernables, cartes bien mélangées -, la situation est équiprobable. La probabilité d’un évènement s’exprime alors par « a chances sur b » : a issues favorables parmi b issues en tout.\end{definition}
\begin{exemple}Obtenir un nombre pair avec un dé : les issues 2, 4, 6 parmi 6, soit 3 chances sur 6 - autant de chances de réussir que d’échouer. Tirer un cœur dans un jeu de 32 cartes : 8 chances sur 32.\end{exemple}
\begin{savoirfaire}[Compter avant de comparer]Comparer deux évènements : compter leurs issues favorables sur le même total. « Plus grand que 4 » (issues 5 et 6) et « au plus 2 » (issues 1 et 2) ont exactement les mêmes chances : 2 sur 6. Le décompte tranche, pas l’intuition.\end{savoirfaire}
\begin{attention}Le dé ne se souvient de rien : après trois 6 de suite, le quatrième lancer garde 1 chance sur 6 de donner un 6. Chaque répétition de l’expérience repart à neuf : c’est l’indépendance des lancers.\end{attention}
\clearpage
\phantomsection\label{cours-21-2}
\essSection{2}{Recenser sans rien oublier}
\begin{definition}Pour une expérience en deux étapes - deux lancers, deux tirages -, un tableau à double entrée ou un arbre recense toutes les issues : chaque case ou chaque chemin est une issue.\end{definition}
\label{fig-cours-21-2-cours-arbre}\figureGeo{cours-arbre}

\begin{exemple}Deux pièces lancées : les issues sont PP, PF, FP, FF, quatre chemins équiprobables. « Obtenir exactement un pile » est réalisé par PF et FP : 2 chances sur 4, une chance sur deux.\end{exemple}
\begin{savoirfaire}[L’arbre, étage par étage]Construire l’arbre étape par étape : les branches de la première étape, puis, au bout de chacune, les branches de la seconde. Le nombre total d’issues est le produit des possibilités : $2\times2=4$ pour deux pièces, $6\times6=36$ pour deux dés.\end{savoirfaire}
\begin{attention}PF et FP sont deux issues différentes : la première pièce et la seconde jouent chacune leur rôle. Les confondre fausse tous les comptes.\end{attention}

\clearpage\rubrique{Découverte}
\begin{blocLivre}{ex-21-1}\begin{exercice}[1]{Premières fractions de chances}
Avec un dé équilibré, donner les chances : d’obtenir 6 ; d’obtenir un nombre pair ; d’obtenir un nombre de 1 à 6.
\end{exercice}\end{blocLivre}
\begin{blocLivre}{ex-21-2}\begin{exercice}[1]{Le sac de boules}
Un sac contient 2 boules rouges, 3 vertes et 5 noires, indiscernables au toucher. Donner les chances de tirer chaque couleur.
\label{fig-ex-21-2}\figureGeo{pro-2}
\end{exercice}\end{blocLivre}
\begin{blocLivre}{ex-21-3}\begin{exercice}[1]{Le jeu de 32 cartes}
Dans un jeu de 32 cartes bien mélangé, donner les chances de tirer : un cœur ; une carte rouge ; l’as de pique.
\end{exercice}\end{blocLivre}
\begin{blocLivre}{ex-21-4}\begin{exercice}[1]{Équiprobable ou pas ?}
La roue A a 8 secteurs identiques ; la roue B a un grand secteur et trois petits. Sur quelle roue peut-on parler de « 1 chance sur 8 » pour chaque secteur ? Pourquoi ?
\label{fig-ex-21-4}\figureGeo{pro-4}
\end{exercice}\end{blocLivre}
\begin{blocLivre}{ex-21-5}\begin{exercice}[1]{Comparer deux évènements}
Avec un dé, quel évènement est le plus probable : « obtenir plus grand que 4 » ou « obtenir au plus 2 » ? Justifier.
\end{exercice}\end{blocLivre}

\rubrique{Entraînement}
\begin{blocLivre}{ex-21-6}\begin{exercice}[2]{La mémoire du dé}
Jade a obtenu trois 6 de suite. Elle affirme : « Le 6 a moins de chances de sortir au prochain lancer, il est fatigué ! » Corriger.
\end{exercice}\end{blocLivre}
\begin{blocLivre}{ex-21-7}\begin{exercice}[2]{Les deux pièces}
On lance deux pièces. Lister toutes les issues dans un arbre, puis donner les chances d’obtenir exactement un pile.
\label{fig-ex-21-7}\figureGeo{pro-7}
\end{exercice}\end{blocLivre}
\begin{blocLivre}{ex-21-8}\begin{exercice}[2]{Le tableau des deux dés}
On lance deux dés et on fait la somme. Combien y a-t-il d’issues en tout ? Combien donnent une somme de 7 ?
\label{fig-ex-21-8}\figureGeo{pro-8}
\end{exercice}\end{blocLivre}
\begin{blocLivre}{ex-21-9}\begin{exercice}[2]{Pile puis dé}
On lance une pièce puis un dé. Combien d’issues compte l’expérience ? Donner les chances d’obtenir « face puis un 6 ».
\end{exercice}\end{blocLivre}
\begin{blocLivre}{ex-21-10}\begin{exercice}[2]{Plus probable que son contraire ?}
Dans un sac, 7 jetons gagnants et 3 perdants. L’évènement « tirer un gagnant » est-il plus probable que son contraire ? Donner les deux comptes de chances.
\end{exercice}\end{blocLivre}

\rubrique{Problèmes}
\begin{blocLivre}{ex-21-11}\begin{exercice}[3]{Les deux tirages de jetons}
Un sac contient un jeton rouge et un bleu. On tire un jeton, on le remet, on retire. Lister les issues, puis donner les chances d’obtenir deux couleurs différentes.
\label{fig-ex-21-11}\figureGeo{pro-11}
\end{exercice}\end{blocLivre}
\begin{blocLivre}{ex-21-12}\begin{exercice}[3]{La loterie honnête ?}
Une tombola vend 200 billets dont 8 gagnants. Maé dit : « J’ai 8 chances sur 200, c’est moins d’une chance sur 20. » A-t-elle raison ?
\end{exercice}\end{blocLivre}
\begin{blocLivre}{ex-21-13}\begin{exercice}[3]{Construire l’expérience}
Inventer une expérience équiprobable où « gagner » a exactement 3 chances sur 8, puis la décrire.
\end{exercice}\end{blocLivre}
\begin{blocLivre}{ex-21-14}\begin{exercice}[3]{Le détour par le contraire}
Avec un dé, donner les chances de « ne pas obtenir 6 », puis expliquer le lien avec les chances d’« obtenir 6 ».
\end{exercice}\end{blocLivre}
\begin{blocLivre}{ex-21-15}\begin{exercice}[3]{Le grand tri}
On tire une carte d’un jeu de 32. Classer du moins probable au plus probable : « tirer l’as de cœur », « tirer un as », « tirer un cœur », « tirer une carte rouge ou noire ».
\end{exercice}\end{blocLivre}

\clearpage
